%% !TEX root = ../manual_gaolv5.tex
% Chapter 5.
%% Chapter 5: Creating intervals. Sources: gaol.tex of f5bf180 lines 2409-2802
%% (interval creation and assignment, assignment combined with an operation,
%% interval constants) and doc/using.md lines 843-898 (integers: the facts
%% about integer bounds).
\chapter{Creating intervals}\label{chap:creation}

An interval of GAOL is an object of the class \code{interval}. A program
creates it from its two bounds, from one number, from nothing, from another
interval, or from the text that writes it; it can also build it from a
midpoint and a radius. Once created, an interval receives new values by
assignment, alone or combined with an operation, and a few intervals of
constant use, the empty set, the whole real line, the enclosures of $\pi$,
are provided ready-made by the class.

\medskip
The way a value enters an interval is not a detail. A \code{double} is taken
as it is, an integer as the integer it is, which may be no double, and a
string as the interval it writes, rounded outward. A decimal constant written
in the program, as $0.1$, is rounded by the compiler before it reaches the
constructor, and an interval built from it does not contain the number the
programmer wrote. The constructors of GAOL v5 answer every input that makes
no interval, bounds in the wrong order, a NaN bound or an infinite point, with
the empty set, and they refuse the silent conversions that GAOL~4 accepted,
from a \code{double} and from a string. Each of these differences is marked
in the text.

\medskip
In this chapter, we describe the creation of intervals and their assignment.
The first section presents the four constructors, the function
\code{textToInterval()}, the static method \code{midrad()} and the intervals
of one infinite point, which GAOL v5 does not have. The second section covers
the straight assignment, the third the assignment combined with an operation
(\code{+=}, \code{\&=}\ldots) and the fourth the interval constants of the
class. The format of the strings that \code{textToInterval()} reads is the
subject of Chapter~\ref{chap:io}.

\section{Constructors}\label{sec:constructors}
\cindex{interval!constructors}%

The class \code{interval} has four constructors. The SSE2 and NEON intervals
also have an explicit constructor from the register that holds the stored
bounds, \code{interval(const \_\_m128d\&)} or
\code{interval(const float64x2\_t\&)}, the lower bound negated in the first
lane, which gives the empty set for a lane of $-\infty$; it serves the code of
GAOL (\pxref{sec:kernels}). A fifth way to create an
interval, from the text that writes it, is the function
\code{textToInterval()}, and a sixth, from a midpoint and a radius, is the
static method \code{interval::midrad()}. We describe the constructors first,
then the two functions, and we end with the intervals of one infinite point.

\subsection{The four constructors}

\begin{itemize}
\item[$\bullet$] \textbf{From the two bounds.} The left and the right bounds
  are given:
  \begin{example}
    interval x(@var^l~,@var^r~);
  \end{example}
  where \var{l} and \var{r} are \code{double}s or of a type that is castable
  into a \code{double}. The result is the empty set when \itv{l}{r} is not an
  interval: a left bound $+\infty$, a right bound $-\infty$, bounds in the
  wrong order ($l>r$) and NaN bounds all give the empty set.\newinvfive{}
  An integer bound of any integer type but \code{bool}, the other bound being
  an integer, a \code{double} or a \code{float}, is taken as the integer it
  is: the left bound is the largest double at most \var{l}, the right one the
  smallest double at least \var{r}. Thus \code{interval(1, 9007199254740993LL)}
  contains $2^{53}+1$, which no double is, and which its conversion to a
  double made $2^{53}$ or $2^{53}+2$ according to the rounding direction. The
  result is the empty set for $l>r$ compared as numbers:
  \code{interval(9007199254740993LL, 9007199254740992LL)} is empty, though
  both bounds lie between the same two doubles.\newinvfive
\cindex{empty!interval: constructing}%

\item[$\bullet$] \textbf{From one number.} A degenerate interval, a point, is
  built from one bound:
  \begin{example}
    interval x(@var^v~);
  \end{example}
  This is equivalent to \code{interval x(\var{v},\var{v});}, and gives the
  empty set for an infinite or NaN \var{v}. The constructor is
  \code{explicit}: a \code{double} no longer converts silently to an
  interval. The statements \code{interval x = 0.5;}, \code{return 0.5;} in a
  function returning an interval, \code{gaol::sin(0.5)} and \code{min(x, 1.0)} do
  not compile, where \code{interval x(0.5)}, \code{interval(0.5)},
  \code{sin(interval(0.5))} and \code{min(x, interval(1.0))} do. The
  assignment (\pxref{sec:assignment}), the operators \code{\&=}, \code{|=}
  and the other assignments combined with an operation
  (\pxref{sec:combined-assignment}), the arithmetic operators and the
  relation symbols (\pxref{sec:relational-symbols}) take a \code{double} as
  they take \code{interval(\var{d})}.\newinvfive{} An integer \var{n} of any
  integer type but \code{bool} gives the tightest interval containing it:
  \itv{n}{n} where $n$ is a double, as every \code{int} is, and the two
  doubles around it otherwise. \code{interval(9007199254740993LL)} is
  \itv{2^{53}}{2^{53}+2}, where the conversion of $2^{53}+1$ to a double gave
  \itv{2^{53}}{2^{53}} or \itv{2^{53}+2}{2^{53}+2}, neither of which contains
  it, and \code{interval(ULLONG\_MAX)} contains $2^{64}-1$. This constructor
  is \code{explicit} too. The same operations take an integer as they take
  \code{interval(\var{n})}; the relations compare it as the integer it is
  (\pxref{sec:relational-symbols}), and \code{\&=} gives
  \code{interval(\var{n})} where the interval contains $n$, the empty set
  otherwise.\newinvfive

\item[$\bullet$] \textbf{Without any bound.} The default constructor builds
  the whole real line:
  \begin{example}
    interval x;
  \end{example}
  This is equivalent to:
  \begin{example}
    interval x(-GAOL_INFINITY,+GAOL_INFINITY);
  \end{example}
  \noindent where \code{GAOL\_INFINITY} is the infinity of the \code{double}
  format (\pxref{sec:floating-point-constants}).

\item[$\bullet$] \textbf{By copy.} The copy constructor duplicates an
  existing interval:
  \begin{example}
    interval x(-12,12), y=x, z(x);
  \end{example}
\end{itemize}

\begin{Rem}
\cindex{integer!as a bound}%
The integers are taken as the integers they are whatever the rounding
direction, and for every integer type but \code{bool}.\newinvfive{} Three
cases still convert to a \code{double}: an unscoped enumeration, an
\code{\_\_int128} outside the GNU modes of GCC and Clang, and a \code{long
double} bound next to an integer, which takes \code{interval(double, double)}.
An integer that is a double gives the result the double gives, and, for the
types of 53 bits or less (\code{int}, \code{short}, \code{char}\ldots),
through the same code. The constructor \code{interval(\var{n})} is
\code{explicit}, as the constructor from a \code{double}. The other
operations of GAOL that take an integer are described in
\exref{sec:integers}.
\end{Rem}

\subsection{Intervals read from a string}

The function \code{textToInterval()} reads a string that writes an interval
in the format of the input (\pxref{sec:input-format}), and returns that
interval. It takes the place of the constructors from strings of GAOL~4.

\begin{deffun}{interval}{textToInterval}{(const std::string\& \var{s})}
  Returns the interval the string \var{s} writes, read with the names of the
  functions of GAOL (\pxref{sec:input-format}). It is in the namespace
  \code{gaol}, under the name IEEE 1788-2015 gives it.
  \begin{example}
interval x = textToInterval("[-23, inf]"),
         y = textToInterval("[5*0.1+dmin, 89*sinh(2.1)]");
  \end{example}
  If the string does not comply with the expected format, the exception
  \code{std::invalid\_argument} is thrown, and \code{std::domain\_error} if it
  calls a function with an argument the function does not take, as
  \code{"nth\_root(8, 1.5)"}, when the library was compiled with exception
  support, which is the default of the three builds. Without exception
  support, the variable \code{errno} is set to \code{-1}, and the function
  \code{gaol\_error()} prints an error message before the program aborts
  (Chapter~\ref{chap:errors}). The \code{textToInterval()} of
  \code{gaol\_ieee1788} reads the names of the functions of the standard, and
  returns the empty set instead (Chapter~\ref{chap:ieee1788}). A null pointer
  is no string: the overload for a \code{const char*}, which
  \code{textToInterval(nullptr)} reaches, throws \code{std::invalid\_argument}
  in both namespaces, and gives the same error without exception support (a
  message, then \code{abort()}).\newinvfive

  GAOL v5 has no constructor from one string. \code{textToInterval()} takes
  the place of the constructor \code{interval(const char*)} of GAOL~4, which
  converted any \code{const char*} implicitly: \code{interval(0)} was
  ambiguous, and \code{x~+~p} compiled for a null \code{p}, and
  crashed.\newinvfive
\end{deffun}

\begin{deffun}{interval}{textToInterval}{(const std::string\& \var{sl}, const std::string\& \var{sr})}
  Returns the interval whose left bound is the one of the interval the string
  \var{sl} writes, and whose right bound is the one of the interval \var{sr}
  writes. The result is the empty set when the lower bound of the first is
  above the upper bound of the second. The comparison is between the
  intervals the two strings write, each rounded outward, not between the
  numbers, where \code{textToInterval("[l,u]")} orders the numbers.
  \begin{example}
 @rem^// x=[-4,4]~
 interval x = textToInterval("[-5,4]+1","[4,6]-[2,3]");
  \end{example}
  It is in the namespace \code{gaol}, and throws \code{std::invalid\_argument}
  as \code{textToInterval(\var{s})} does, for a null pointer in either place
  too. It takes the place of the constructor from two strings of GAOL~4, which
  kept the bounds \itv{2}{1} of \code{interval("2",~"1")}: \code{is\_empty()}
  was true, but its sum with \itv{0}{1} was \itv{2}{2}.\newinvfive
\end{deffun}

\subsection{A midpoint and a radius}

The static method \code{interval::midrad()} builds an interval from its
midpoint and its radius, both given as \code{double}s.

\begin{defstaticmethod}{interval}{interval}{midrad}{(double \var{m}, double \var{r})}
  \begin{operation}
    \op{$\code{midrad}(m,r)=\itv{\roundDn{m-r}}{\roundUp{m+r}}$}
  \end{operation}
  Returns the interval of midpoint \var{m} and radius \var{r}, computed as
  \code{interval(\var{m}) + interval(-\var{r}, \var{r})}. It encloses every
  real of \itv{m-r}{m+r}, where \code{interval(\var{m} - \var{r}, \var{m} +
  \var{r})} rounds both bounds in the same direction, upward unless GAOL is built
  with \option{GAOL\_PRESERVE\_ROUNDING}. A radius $r=0$ gives the point \var{m},
  and $r=+\infty$ the whole line. A negative or NaN \var{r} gives the empty
  set, and so does an infinite or NaN \var{m}, as \code{interval(\var{m})}
  does.\newinvfive
  \begin{example}
cout << interval::midrad(1, 0.25) << ' '
     << interval::midrad(1, -1);
@outputs^[0.75, 1.25] [empty]~
  \end{example}
\end{defstaticmethod}

\subsection{Decimal constants}

Of the three forms in which a number reaches an interval, the \code{double} is
the usual one, and the one that calls for care.

\begin{Rem}
A rational number with a finite expansion in the decimal base may need a
rounding in the binary base. The statement
\begin{example}
  interval x(0.1);
\end{example}
\noindent
does \emph{not} create an interval containing $0.1$. This number has an
infinite expansion in base two, so no significand, whatever its size,
represents it exactly, and the constant $0.1$ is rounded, in all likelihood
to the nearest floating-point number, when the program is compiled
(\pxref{sec:rounding-definitions}, and \exref{sec:common-errors}). The
interval that contains $0.1$ is read from a string:
\begin{example}
interval x = textToInterval("0.1");
\end{example}
\end{Rem}

\subsection{Infinite point intervals}\label{par:inf_inf_note}
\cindex{interval!with infinite bounds}%

The degenerate intervals \itv{+\infty}{+\infty} and \itv{-\infty}{-\infty} do
not exist in GAOL v5.\newinvfive{} IEEE 1788-2015 \cite{IEEE1788:2015} has no
such interval: its constructor fails there and gives the empty set (clauses
10.5.8 and 12.12.7). The constructors of GAOL give the empty set too, as
those of IBEX do, for an infinite point and for all the cases listed above:
\begin{example}
cout << interval(GAOL_INFINITY) << ' '
     << interval(-GAOL_INFINITY,-GAOL_INFINITY) << ' '
     << interval(2,1) << ' '
     << interval(1,GAOL_INFINITY) << endl;
@outputs^[empty] [empty] [empty] [1, inf]~
\end{example}
GAOL~4 built them as intervals of one infinite point, which its operators did
not support. The intervals with one infinite bound, as \itv{1}{+\infty}, are
intervals as any other.

\medskip
Read from a string, the constant \element{inf}, an expression, stands for
\itv{\text{MAX}}{+\infty}, MAX being the largest double, and \element{-inf}
for \itv{-\infty}{-\text{MAX}}. The interval literals whose left bound is
$+\infty$ or whose right bound is $-\infty$ are the empty set, as
\ieee{numsToInterval} of IEEE 1788-2015 has it:
\begin{example}
cout << textToInterval("inf") << ' ' << textToInterval("[inf]") << ' '
     << textToInterval("[-inf, -inf]") << endl;
@outputs^[1.797693134862315e+308, inf] [empty] [empty]~
\end{example}
The string \element{"<-inf, -inf>"}, which GAOL~4 read as
\itv{-\infty}{-\text{MAX}}, throws \code{std::invalid\_argument}: a bound of
\code{<a, b>} must have a double for its value, and the value of \element{-inf}
is \itv{-\infty}{-\text{MAX}}, no point (\pxref{sec:input-format}).\newinvfive

\section{Straight assignment}\label{sec:assignment}
\cindex{assignment}%

An existing interval receives a new value from another interval, from a string,
from a \code{double}, or from an integer (below).

\begin{itemize}
\item[$\bullet$] \textbf{From another interval.} The assignment copies it:
  \begin{example}
    interval x(-12,12),
             y; @rem^// Here, y is @itv^-@infty~^+@infty~~
    y = x; @rem^// Now, y is @itv^-12~^12~~
  \end{example}
\item[$\bullet$] \textbf{From a string.} The string follows the format
  expected for the input (\pxref{sec:reading-intervals}), and
  \code{textToInterval()} reads it:

\begin{example}
interval x;
x = textToInterval("[-inf, 123]");
\end{example}
  GAOL~4 assigned the string itself, \code{x~=~"[-inf,~123]"}. GAOL v5
  refuses this statement, having no constructor from one string.\newinvfive

\item[$\bullet$] \textbf{From a \code{double}.} The operator
  \code{operator=(double)} assigns it as \code{interval(\var{d})}, the
  constructor from one \code{double} being explicit:\newinvfive

\begin{example}
interval x;
x = 1234.5; @rem^// Now, x is @itv^1234.5~^1234.5~~
x = GAOL_INFINITY; @rem^// Now, x is empty~
\end{example}
\end{itemize}

\noindent
The assignments from a string and from a \code{double} follow the rules of
\exref{sec:constructors}: assigning an infinity to an interval empties it. An
integer \var{n} is assigned as \code{interval(\var{n})}, the tightest interval
containing it.\newinvfive

\begin{Rem}
The interface has no method for modifying one bound of an interval: an interval
is created, assigned and replaced as a whole. (The non-const accessors
\code{get\_xmminterval()} and \code{get\_neoninterval()} of the SSE2 and NEON
intervals hand the register of the stored bounds to the code of GAOL; a program
does not write it.)
\end{Rem}

\section{Assignment combined with an operation}\label{sec:combined-assignment}

In this manual, we write \code{self} for the object to which \code{this} is
a pointer in C++.\cindex{this@\texttt{this}}\cindex{self@\texttt{self}} The
assignment operators below combine the value of \code{self} with the value
of their right-hand side, an interval, a \code{double} or an integer, and store the
result in \code{self}.

\defmethodx{interval}{interval\&}{operator\&=}{(const interval\& \var{I})}
\begin{defmethod}{interval}{interval\&}{operator\&=}{(double \var{d})}
  \begin{operation}
    \op{$\code{self}\gets\code{self} \cap I$}
    \op{$\code{self}\gets\code{self} \cap \code{interval}(d)$}
  \end{operation}
Assigns to \code{self} the interval resulting from the intersection of
\code{self} and \var{I} (resp. \code{interval(\var{d})}, empty for an
infinite or NaN \var{d}), which is the empty set when they are
disjoint.\newinvfive

\begin{example}
interval x(-12,12);
x &= interval(-6,23); @rem^// Now, @code^x~ is @itv^-6~^12~~
\end{example}

\end{defmethod}

\defmethodx{interval}{interval\&}{operator\textbar=}{(const interval\& \var{I})}
\begin{defmethod}{interval}{interval\&}{operator\textbar=}{(double \var{d})}
  \begin{operation}
    \op{$\code{self}\gets\hull{(\code{self} \cup I)}$}
    \op{$\code{self}\gets\hull{(\code{self} \cup \code{interval}(d))}$}
  \end{operation}
  Assigns to \code{self} the interval hull of the union of \code{self} and
  \var{I} (resp. \code{interval(\var{d})}, which leaves \code{self} as it is
  for an infinite or NaN \var{d}).\newinvfive

\begin{example}
interval x(-12,12);
x |= interval(-6,23); @rem^// Now, @code^x~ is @itv^-12~^23~~
\end{example}

\end{defmethod}

\defmethodx{interval}{interval\&}{operator+=}{(const interval\& \var{I})}
\begin{defmethod}{interval}{interval\&}{operator+=}{(double \var{d})}
  \begin{operation}
    \op{$\code{self} \gets \code{self} + I$}
    \op{$\code{self} \gets \code{self} + d$}
  \end{operation}

  Assigns to \code{self} the interval resulting from adding \code{self} and
  \var{I} (resp. \var{d}).

\begin{example}
interval x(-12,12);
x += interval(-6,23); @rem^// Now, @code^x~ is @itv^-18~^35~~
\end{example}
\end{defmethod}

\defmethodx{interval}{interval\&}{operator-=}{(const interval\& \var{I})}
\begin{defmethod}{interval}{interval\&}{operator-=}{(double \var{d})}
  \begin{operation}
    \op{$\code{self}\gets \code{self} - \var{I}$}
    \op{$\code{self}\gets \code{self} - \var{d}$}
  \end{operation}
  Assigns to \code{self} the interval resulting from subtracting \var{I}
  (resp. \var{d}) from \code{self}.
\begin{example}
interval x(-12,12);
x -= interval(-6,23); @rem^// Now, @code^x~ is @itv^-35~^18~~
\end{example}
\end{defmethod}

\defmethodx{interval}{interval\&}{operator*=}{(const interval\& \var{I})}
\begin{defmethod}{interval}{interval\&}{operator*=}{(double \var{d})}
\begin{operation}
\op{$\code{self}\gets \code{self} \times\var{I}$}
\op{$\code{self}\gets \code{self} \times\var{d}$}
\end{operation}
Assigns to \code{self} the interval resulting from multiplying \code{self}
and \var{I} (resp. \var{d}).
\begin{example}
interval x(-12,12);
x *= interval(-6,23); @rem^// Now, @code^x~ is @itv^-276~^276~~
\end{example}
\end{defmethod}

\defmethodx{interval}{interval\&}{operator/=}{(const interval\& \var{I})}
\begin{defmethod}{interval}{interval\&}{operator/=}{(double \var{d})}
\begin{operation}
\op{$\code{self}\gets \code{self} / \var{I}$}
\op{$\code{self}\gets \code{self} / \var{d}$}
\end{operation}
Assigns to \code{self} the interval resulting from dividing \code{self} by
\var{I} (resp. \var{d}).
\begin{example}
interval x(-12,12);
x /= interval(-6,23);   @rem^// Now, @code^x~ is @itv^-@infty~^+@infty~~
x /= interval::zero();  @rem^// Now, @code^x~ is @ensuremath^@emptyset~~
\end{example}
\end{defmethod}

\defmethodx{interval}{interval\&}{operator\%=}{(const interval\& \var{I})}
\begin{defmethod}{interval}{interval\&}{operator\%=}{(double \var{d})}
Assigns to \code{self} the interval resulting from dividing \code{self} by
\var{I} (resp. \var{d}), using a \emph{relational division}
(\pxref{sec:relational-arithmetic}).
\begin{example}
interval x(-12,12);
x %= interval(-6,23);   @rem^// Now, @code^x~ is @itv^-@infty~^+@infty~~
x %= interval::zero();  @rem^// Now, @code^x~ is @itv^-@infty~^+@infty~~
\end{example}
\end{defmethod}

\noindent
The arithmetic assignments with a \code{double} (\code{+=}, \code{-=},
\code{*=}, \code{/=}, \code{\%=}) give the result of the same operation with
\code{interval(\var{d})}, choosing the case by the sign of \var{d}; an infinite
or NaN \var{d} gives the empty set, as \code{interval(\var{d})} is empty.\newinvfive

\section{Interval constants}\label{sec:interval-constants}
\cindex{interval!constants}%

Some useful intervals, and some canonical intervals enclosing real constants,
are defined as static methods of the class \code{interval}.
Table~\ref{tab:interval-constants} lists them.

\begin{table}[htbp]
  \centering\renewcommand{\arraystretch}{1.1}
  \caption{The interval constants, static methods of the class \protect\code{interval}} \label{tab:interval-constants}
  \begin{tabular}{lc}
    \hline
    \multicolumn{1}{c}{Function} & \multicolumn{1}{c}{Value}\\
    \hline\hline
    \code{interval::emptyset()} & $\emptyset$\\
    \code{interval::half\_pi()} &
      \itv{\roundDn{\frac{\pi}{2}}}{\roundUp{\frac{\pi}{2}}}\\
    \code{interval::minus\_one\_plus\_one()} & \itv{-1}{1}\\
    \code{interval::negative()} & \itv{-\infty}{0}\\
    \code{interval::one()} & $1$\\
    \code{interval::one\_plus\_infinity()} & \itv{1}{+\infty}\\
    \code{interval::pi()} & \itv{\roundDn{\pi}}{\roundUp{\pi}}\\
    \code{interval::positive()} & \itv{0}{+\infty}\\
    \code{interval::two\_pi()} & \itv{\roundDn{2\pi}}{\roundUp{2\pi}}\\
    \code{interval::universe()} & \itv{-\infty}{+\infty}\\
    \code{interval::zero()} & $0$\\
    \hline
  \end{tabular}

\end{table}

\noindent
The enclosures of $\pi$, $2\pi$ and $\pi/2$ are the tightest ones
(\pxref{sec:accuracy-modes}). The empty set is held as an interval whose two
bounds are NaN.
\cindex{empty!interval: representation}%

\begin{example}
cout << interval::emptyset() << ' ' << interval::pi() << endl;
@outputs^[empty] [3.141592653589793, 3.141592653589794]~
\end{example}

\section{Discussion}

This chapter described the ways an interval of GAOL comes to exist and to
change: four constructors, the function \code{textToInterval()} for the text
of an interval, the static method \code{midrad()} for a midpoint and a
radius, the straight assignment, the assignments combined with an operation,
and the constants of the class. One rule runs through the constructors,
\code{midrad()} and the assignments: an input that makes no interval, bounds in
the wrong order, a NaN bound or an infinite point, gives the empty set rather
than an object the operators do not support; \code{textToInterval()} gives it
too for a literal such as \code{"[2, 1]"} or \code{"[inf]"}, but throws for a
string that is no expression, \code{"<-inf, -inf>"} among them. GAOL v5 follows IEEE 1788-2015 and IBEX there, where GAOL~4 built
intervals of one infinite point.

\medskip
Two choices of GAOL v5 cost the programmer a few characters and remove
ambiguities. The constructors from one \code{double} and from one integer are
explicit, so that a number no longer becomes an interval without being asked
to. The constructors from strings of GAOL~4 are replaced by
\code{textToInterval()}, so that \code{interval(0)} is no longer ambiguous
and a null pointer no longer compiles into a crash. An integer
bound is taken as the integer it is, whatever the rounding direction, which
keeps $2^{53}+1$ inside the interval that encloses it. The decimal constants
remain the responsibility of the programmer: the way to an interval
containing $0.1$ is a string.

\medskip
The intervals created here are compared in the next chapter, which describes
the set relations, the certainly relations and the order of IEEE 1788-2015 on
intervals (Chapter~\ref{chap:relations}).
